% Gospel-only replacement fragments.  The canonical leaf invokes these macros
% at the page-two dossier and detailed-commentary positions.
\newcommand{\TenthGospelDossier}{%
\textbf{Gospel} & Luke 18:9--14 & Provenance uncertain & commonly c.~AD 80--90\\
\cmidrule(lr){1-4}
\emph{Parable setting} & Luke 18:9--14 & Jerusalem Temple & Jesus' final journey\\
\cmidrule(lr){1-4}
\multicolumn{4}{>{\raggedright\arraybackslash}p{0.92\linewidth}}{Ancient tradition identifies Luke, the Antiochene physician and Pauline companion, and later tradition locates his writing in Achaia. Modern Catholic orientation commonly places the Gospel's final form about AD~80--90, for a substantially Gentile Christian audience outside Palestine; the city is unknown. The Temple is the setting inside Jesus' parable. Luke does not report that two independently identifiable men prayed there on a datable day, so the setting must not be turned into a historical event at c.~AD~29--30 (Luke introduction and 18, \emph{NABRE}).}\\[0.1em]
}

\newcommand{\TenthGospelDetailedCommentary}{%
\sectionguard
\subsection*{The publican's petition and the Pharisee's comparison\enspace\properrefs{Gosp.}}

Luke supplies his interpretive control before the story: Jesus addresses
``some'' who trusted in themselves that they were just and despised the rest.
The two men and the Temple belong to the parable's narrated world; they are
not two independently recoverable worshipers or a datable Temple incident.
The narrative follows the persistent widow and precedes Jesus' reception of
children, the rich ruler, and the blind beggar. It therefore tests prayer,
status, dependence, and the reception of God's kingdom across a larger
Lukan sequence, not merely two bodily postures.

The Jewish setting needs equal care. Fasting and tithing are not condemned:
the Pharisee's fault is making real observances into a comparative claim
against another person. ``Pharisee'' cannot serve as a cipher for Judaism or
the Jewish people, and Luke's own Gospel presents Pharisees in more than one
role. Nor is the publican's occupation endorsed. Tax collectors worked within
an imperial revenue order vulnerable to abuse; elsewhere Luke requires
repentance and restitution from Zacchaeus. The reversal concerns self-trust,
contempt, confession, and God's judgment, not an ethnic contrast or a
declaration that vice is preferable to observance.

Three Greek details sharpen rather than settle the reading. The phrase
\latin{pros heauton} can be construed with the Pharisee's standing or with
his praying, so it should not bear a complete theology by punctuation alone.
The publican's \latin{hilastheti moi} asks God to deal mercifully or
propitiatingly with ``me, the sinner''; the passive form keeps the saving
action with God. The comparative phrase \latin{par' ekeinon} has generated
discussion over whether the publican goes home justified ``rather than'' or
``in comparison with'' the other. Either construction remains governed by
Luke's explicit conclusion: everyone exalting himself will be humbled, and
the one humbling himself will be exalted.

Tertullian reads the scene against Marcion to insist that the same Creator
who condemns pride justifies the humble petitioner
(\emph{Adversus Marcionem} IV.36). Augustine,
\emph{Sermon 115.2--3} (the same sermon numbered 65 in an older English
series), distinguishes the Pharisee's self-praise and accusation from the
publican's confession and pardon; his anti-Pelagian application is reception,
not a warrant that spoken self-accusation works automatically. Cyril's
\emph{Sermon 120 on Luke} likewise contrasts boastful comparison with
embodied confession and leaves justification as God's act.

Later reception develops different consequences without collapsing them.
Chrysostom uses the parable at \emph{On the Statues} 3.8 to show that fasting
without the accompanying life of virtue is fruitless, and at 3.13 to show
that even a true charge does not license detraction. The \emph{Rule of
Benedict} 7 places the publican's lowered gaze and breast-beating near the
summit of humility, but does not make posture self-validating. Anthony of
Padua's historically numbered Eleventh Sunday after Pentecost interprets
Luke 18 with a different Introit and Epistle, so it witnesses medieval
reception of the Gospel rather than this complete 1962 formulary.

The Byzantine Sunday of the Publican and Pharisee gives the parable a distinct
liturgical afterlife at the opening of the Triodion and preparation for Lent.
Its liturgical use praises humility while refusing the inference that good
works themselves are evil. This reception neither derives from nor explains
the Roman Sunday. No checked source established the familiar Jesus Prayer as
a simple historical development from Luke 18:13; the verbal affinity may be
noted only as affinity, not origin.
}
